% explainer-source-hash: 52c32d41dd9ee9308816b0ea13c8f6fd4352d582538ab1ea52c82ccf33c17cf5
% explainer-source-commit: 4226652b05a73b52389e480fbcab2170bbee92a3
% explainer-generated: 2026-07-16
% Regenerate with /explain-all (see WIRING.md). Do not edit this stamp by hand:
% it is how the toolchain knows whether this document still matches the code.
\documentclass[11pt,a4paper]{article}
\input{preamble}

\title{\vspace{-1.5cm}\textbf{Breakout Game}\\[0.3cm]
\large The Brief: What It Is, How It Works, What Is Wrong With It}
\author{Portfolio explainer: high-level summary}
\date{}

\begin{document}
\maketitle
\thispagestyle{fancy}

\section{The project in one paragraph}

\file{breakout-game} is a clone of the 1976 arcade game Breakout: a paddle, a bouncing
ball, a wall of coloured blocks, three lives, and a persistent high score. It is 475 lines
of Python across five modules and it uses \textbf{nothing but the standard library}. The
point of the project is what it refuses to use: there is no game engine, so the frame loop,
the collision detection, the state machine and the difficulty curve are all written by hand
on top of \code{turtle}, a children's drawing module that was never meant to carry a game.

\begin{keyidea}{Why that constraint is the whole story}
An engine like Pygame or Godot hands you a render loop, sprites and collision tests already
solved. Removing it means writing each of them yourself, which is the exercise. Everything
interesting in this codebase, and every bug in it, lives in the parts an engine would have
hidden.
\end{keyidea}

\section{Architecture}

\file{main.py} owns the loop and the game state. The other four modules are passive: they
only act when the loop calls a method on them, and none of them knows a frame exists.

\begin{figure}[H]
\centering
\resizebox{0.96\textwidth}{!}{
\begin{tikzpicture}
  \node[box, minimum width=50mm] (main) {\textbf{main.py} \ (134)\\Screen, state machine, frame loop};
  \node[extbox, above=12mm of main] (turtle) {Python stdlib: \code{turtle} (on \code{tkinter})};

  \node[box, minimum width=30mm] at ([xshift=-65mm]$(main.south)+(0,-30mm)$) (paddle)
    {\textbf{paddle.py} \ (34)\\keys, clamping};
  \node[box, minimum width=30mm] at ([xshift=-21mm]$(main.south)+(0,-30mm)$) (ball)
    {\textbf{ball.py} \ (102)\\velocity, wall + paddle\\collisions, difficulty};
  \node[box, minimum width=30mm] at ([xshift=21mm]$(main.south)+(0,-30mm)$) (blocks)
    {\textbf{blocks.py} \ (147)\\wall layout, block\\collisions, effects};
  \node[box, minimum width=30mm] at ([xshift=65mm]$(main.south)+(0,-30mm)$) (score)
    {\textbf{scoreboard.py} \ (58)\\score, lives};
  \node[store, below=14mm of score] (file) {highestScore\\.txt};

  \draw[dflow] (turtle) -- (main);
  \draw[flow] (main.south) -- (ball.north);
  \draw[flow] (main.south) -- (paddle.north);
  \draw[flow] (main.south) -- (blocks.north);
  \draw[flow] (main.south) -- (score.north);
  \draw[flow] (blocks.south) to[out=-90,in=-90] node[lbl,below]{\code{.increase\_score()}} (score.south);
  \draw[flow] (score.east) to[out=0,in=0] node[lbl,right]{read at import,\\write at game end} (file.east);
  \draw[flow] (ball.south) to[out=-90,in=-90] node[lbl,below]{reads \code{paddle.xcor()}} (paddle.south);
\end{tikzpicture}}
\caption{Five modules, line counts measured. The split is by \emph{collision problem}: ball
versus wall, ball versus paddle and ball versus block are three different shapes of maths,
and each lives with the object that owns it.}
\end{figure}

\section{How a frame works}

\begin{figure}[H]
\centering
\resizebox{0.88\textwidth}{!}{
\begin{tikzpicture}[node distance=6mm,
  flowstep/.style={rectangle, rounded corners=2pt, draw=inkblue, fill=softgrey,
                   text=inkblue, align=center, inner sep=4pt, minimum height=7mm,
                   minimum width=52mm, font=\small}]
  \node[flowstep, fill=white, draw=linegrey] (top) {\code{while True:} \ \ \ \ \code{wn.update()}};
  \node[dec, below=8mm of top, minimum width=20mm] (st) {state ==\\``playing''?};
  \node[flowstep, below=10mm of st] (phys) {\textbf{move} the ball \ \code{(x+=dx, y+=dy)}};
  \node[flowstep, below=of phys] (coll) {\textbf{collide}: blocks, walls, paddle, blocks again};
  \node[flowstep, below=of coll] (fx) {\textbf{advance} flash + particle effects};
  \node[dec, below=9mm of fx, minimum width=22mm] (end) {wall cleared?\\ball below -280?};
  \node[flowstep, below=9mm of end, draw=warnred] (res) {save high score, show end screen,\\or lose a life and re-centre};
  \node[flowstep, left=22mm of st, draw=linegrey] (idle) {effects only,\\no physics};

  \draw[flow] (top) -- (st);
  \draw[flow] (st) -- node[lbl,right]{yes} (phys);
  \draw[flow] (st) -- node[lbl,above]{no} (idle);
  \draw[flow] (phys) -- (coll); \draw[flow] (coll) -- (fx); \draw[flow] (fx) -- (end);
  \draw[flow] (end) -- node[lbl,right]{yes} (res);
  \draw[flow] (end.west) -- node[lbl,above]{no} ++(-3.6,0) |- (idle.south);
  \draw[flow] (idle.north) -- ++(0,0.5) -| (top.west);
\end{tikzpicture}}
\caption{One frame of \file{main.py}. Physics runs only in the \code{playing} state; the
other states still spin the loop so the window stays responsive to a keypress.}
\end{figure}

\begin{keyidea}{The one line that makes it a game}
\code{wn.tracer(0)} switches off turtle's redraw-after-every-command behaviour, and the loop
calls \code{wn.update()} exactly once per pass. That converts turtle from ``draw each change
instantly and slowly'' into proper frame-based rendering: move everything invisibly, then
present one complete image. It is the single non-obvious idea in the project.
\end{keyidea}

\subsection{State machine}

\begin{figure}[H]
\centering
\begin{tikzpicture}[node distance=20mm]
  \node[box] (start) {\textbf{start}};
  \node[box, right=of start] (playing) {\textbf{playing}};
  \node[box, above right=10mm and 16mm of playing] (won) {\textbf{won}};
  \node[box, below right=10mm and 16mm of playing] (over) {\textbf{game\_over}};
  \draw[flow] (start) -- node[lbl,above]{\code{space}} (playing);
  \draw[flow] (playing) -- node[lbl,above left=-2mm and -3mm]{wall cleared} (won);
  \draw[flow] (playing) -- node[lbl,below left=-2mm and -3mm]{last life lost} (over);
  \draw[flow] (won.north) to[out=100,in=60, looseness=1.4] node[lbl,above]{\code{r}} (start.north);
  \draw[flow] (over.south) to[out=-100,in=-60, looseness=1.2] node[lbl,below]{\code{r}} (start.south);
\end{tikzpicture}
\caption{Four states held in one string variable. Both key callbacks are guarded, so
\code{r} is a no-op mid-rally and \code{space} is a no-op on the end screen.}
\end{figure}

\section{The three collision rules}

\begin{center}
\begin{tabular}{@{}p{0.19\textwidth}p{0.31\textwidth}p{0.42\textwidth}@{}}
\toprule
\textbf{Collision} & \textbf{Test} & \textbf{Response} \\
\midrule
Ball vs.\ wall & \code{|x| > 390} or \code{|y| > 290} & clamp back onto the boundary, then
flip that axis. The clamp matters: flipping alone leaves the ball outside the wall, where it
vibrates forever. \\
\addlinespace
Ball vs.\ paddle & \code{y < -240} \emph{and} within 80 units of the paddle's centre &
clamp to \code{y = -240}, flip \code{dy}, count the hit. Directional on purpose: a ball
moving upward passes through, which is what stops it sticking to the paddle. \\
\addlinespace
Ball vs.\ block & \code{block.distance(ball) < 35} & score, remove, flash white, burst, flip
\code{dy}. Turtle has no rectangle intersection test, so a fixed radius stands in for every
block's true hitbox. \\
\bottomrule
\end{tabular}
\end{center}

\section{The rest of the design}

\begin{itemize}
  \item \textbf{Wall generation} (\file{blocks.py}) packs randomly-sized blocks right to
        left from $x = 400$, wraps down 20 units when the next will not fit, and stops when
        it runs out of vertical room. Each row's colour comes from its $y$ position, so the
        wall reads as bands rather than a speckle.
  \item \textbf{Difficulty ramp}: every 4th paddle return multiplies both velocity
        components by 1.08. Multiplying rather than adding preserves direction, which is why
        it is written that way.
  \item \textbf{Two resets}, and the distinction is a real design decision: losing a life
        keeps the current speed (a setback), while restarting after game over drops back to
        base speed.
  \item \textbf{Break feedback}: a broken block flashes white for 3 frames while five
        particle turtles fly outward on random angles, driven by a frame countdown rather
        than any physics.
  \item \textbf{Persistence}: the high score is read from \file{highestScore.txt} when
        \file{scoreboard.py} is imported, and written whenever a game ends.
\end{itemize}

\section{Honest assessment}

Every finding below was confirmed by \textbf{running} the project's real modules, not by
reading them. The Deep Dive has the evidence and the traces.

\begin{honest}{The five real defects}
\begin{enumerate}
  \item \textbf{A ghost block scores a free point on frame one of every game.} When the wall
        runs out of room, the last block is left at its default position, the origin, and
        appended to \code{all\_blocks} anyway; \code{make\_all\_blocks()} hides it but never
        removes it. The ball also starts at the origin, so it collides immediately: the
        player gets a free point, the ball launches \emph{downward} instead of up, and a
        steel-blue particle burst fires from the middle of the screen. Every symptom looks
        like a feature if you have not read the code.
  \item \textbf{A first run starts the high score at 1.} \code{score = open(...).write("0")}
        assigns \code{write()}'s \emph{return value}, which is the number of characters
        written. The README claims this branch starts the score at 0.
  \item \textbf{The bottom wall is unreachable.} It bounces at \code{y < -290}, but the
        life-lost check fires at \code{y < -280} and the ball moves at most 2.6 units per
        frame. The README says the ball bounces off all four walls; it bounces off three.
  \item \textbf{The double collision check does nothing.} \file{main.py} calls
        \code{check\_collision\_ball()} twice per frame, apparently against tunnelling, but
        both calls test position rather than the swept path, and a 2.6-unit step against a
        35-unit radius cannot tunnel anyway. It can, however, destroy two blocks in one frame
        and flip \code{dy} twice, cancelling the bounce out.
  \item \textbf{Two of the seven row colours are unreachable.} The geometry allows exactly
        five rows, so \code{'medium purple'} never appears and \code{'steel blue'} only ever
        colours the ghost block. The comment about cycling colours describes a case that
        cannot happen.
\end{enumerate}
\end{honest}

\begin{honest}{The design gap that matters most}
The paddle is a perfect mirror: \code{dy} flips and \code{dx} is never touched, so
\emph{where} the ball hits the paddle has no effect. In real Breakout, angling the ball off
the paddle's edge is the entire skill of the game. Here the player cannot aim, which means
there is no strategy to play, only survival. Deriving \code{dx} from the strike offset is a
handful of lines and is the highest-value change available to the project.

Related, and cheaper: there is no frame-rate limiter and no delta-time scaling, so the game
literally runs faster on a faster computer, and the tuned constants mean nothing in absolute
terms.
\end{honest}

\begin{honest}{Documentation drift}
The README quotes ``279 lines'' (measured: 475), says the game imports \code{time} (nothing
does; the real imports are \code{turtle}, \code{random} and \code{math}), and describes the
missing-file high score as 0 (it is 1). The README's judgement is sound, and its ``What I'd
do differently'' independently identifies three of the genuine weaknesses. The drift is in
the specifics.
\end{honest}

\subsection{What it gets right}

\begin{itemize}
  \item The module split earns itself: each collision rule is a different problem and lives
        with the object that owns it.
  \item \code{tracer(0)} plus a manual \code{update()} is the correct, non-obvious way to
        get frame-based rendering out of turtle.
  \item Clamp-then-flip shows the author hit the stuck-ball bug and fixed it properly, and
        the directional paddle test shows the same about the sticky-paddle bug.
  \item \code{for entry in self.fading[:]} shows real awareness of the
        mutation-during-iteration hazard, though \code{check\_collision\_ball()} only avoids
        the same hazard by returning early, which is correct by accident.
  \item Catching \code{Terminator} so closing the window exits cleanly, rather than dumping
        a traceback, is a detail most projects this size skip.
\end{itemize}

\begin{keyidea}{The summary}
A competent exercise that does what it set out to do: a working game loop, a state machine
and three kinds of hand-rolled collision detection, with a file structure the author can
justify. It also ships an invisible block that scores a free point, a high score that starts
at 1, three walls where the README promises four, and a paddle that cannot be aimed with.
All are small fixes, and each is a better story told first than discovered by an interviewer.
\end{keyidea}

\section{Glossary}

\begin{description}[leftmargin=0pt, style=nextline, itemsep=3pt]
  \item[Game engine] A library that solves the hard parts of a game for you: render loop,
        sprites, collision, input, physics. This project deliberately uses none.
  \item[\code{turtle}] Python's standard-library drawing module, built on \code{tkinter}. A
        cursor you move around a window. Not a game library: no sprites, no collision
        detection, no frame timing.
  \item[Frame] One still image of the game. A game draws one, moves everything slightly,
        draws the next, forever.
  \item[Game loop] The \code{while True} that runs until the player quits. Each pass is one
        frame: move, collide, draw.
  \item[Velocity / \code{dx}, \code{dy}] Speed plus direction, stored as horizontal and
        vertical change per frame. Moving is addition; bouncing is multiplying one of them
        by $-1$.
  \item[Hitbox] An invisible rectangle used for collision maths instead of an object's drawn
        pixels. This project has none: it substitutes a fixed-radius circle for every block.
  \item[Tunnelling] When a fast object jumps clean over a thin one between two frames, so a
        position-based collision test never sees the overlap. Impossible at this project's
        speeds, which is why the defence against it is wasted work.
  \item[State machine] A design where the program is always in exactly one named state and
        moves between them on defined events. Here: one string, four states, two key events.
  \item[Persistence] Keeping data after the process exits, by writing it to a file. The one
        persisted value here is the high score.
  \item[\code{Terminator}] The exception \code{turtle} raises when its window is closed.
        Caught in \file{main.py} so the process exits quietly.
\end{description}

\end{document}
